% CHƯƠNG 3. DỮ LIỆU VÀ PHƯƠNG PHÁP
\clearpage
\section[Dữ liệu và phương pháp]{DỮ LIỆU VÀ PHƯƠNG PHÁP}

Chương này giới thiệu sáu bộ dữ liệu (mục 3.1), cách xử lý và khám phá dữ liệu (mục 3.2), các phương pháp kiểm chứng dùng chung cho mọi nội dung (mục 3.3) và thang xếp loại kết quả (mục 3.4). Các chương 4, 5, 6 chỉ nhắc lại những điểm riêng của từng nội dung.

\subsection{Mô tả dữ liệu}

Cả sáu bộ dữ liệu lấy từ kho Rdatasets, bản sao các bộ dữ liệu của gói \texttt{datasets} trong R [DL1], tải ngày 20/11/2026. Chúng chia làm hai nhóm: Hồ Huron, Linh miêu và DAX (suất sinh lợi) cho các mô hình dừng ở Chương 4; Nile, AirPassengers và DAX (log giá) cho chuỗi không dừng và dự báo ở Chương 5 và 6 (Bảng~\ref{tab:dulieu}).

Chỉ số DAX được dùng ở cả hai nhóm dưới hai dạng. Log giá đóng cửa trôi theo thời gian và không dừng; suất sinh lợi log, tức sai phân của log giá nhân với 100, là dạng dừng của chính chuỗi đó. Dùng cùng một chuỗi gốc cho hai phía giúp thấy rõ sai phân biến chuỗi không dừng thành dừng. Dữ liệu DAX trong gói R lấy theo ``thời gian kinh doanh'' (bỏ cuối tuần và ngày lễ) và không có ngày lịch, nên trục thời gian là số thứ tự ngày giao dịch.

\begin{table}[H]
\centering\footnotesize
\caption{Sáu bộ dữ liệu được sử dụng}\label{tab:dulieu}
\begin{adjustbox}{max width=\linewidth}
\begin{tabular}{@{}p{2.6cm}p{2.8cm}p{2.0cm}p{4.0cm}rp{2.6cm}p{3.2cm}@{}}
\toprule
Bộ dữ liệu & Nguồn (R) & Quan sát & Biến được mô hình hóa & $n$ & Thời gian & Phép biến đổi \\
\midrule
Hồ Huron & \texttt{LakeHuron} [DL2] & hằng năm & mực nước hồ (ft) & 98 & 1875--1972 & không; biến xu thế $x=\text{năm}-1920$ \\
Linh miêu & \texttt{lynx} [DL3] & hằng năm & số linh miêu bị bẫy ở Canada & 114 & 1821--1934 & $\log_{10}$ \\
DAX (suất sinh lợi) & \texttt{EuStockMarkets}, cột DAX [DL4] & ngày giao dịch & suất sinh lợi log (\%) & 1859 & $t=2,\dots,1860$ & $100\,\Delta\log P_t$ \\
Nile & \texttt{Nile} [DL5] & hằng năm & lưu lượng sông Nile tại Aswan ($10^8\,\mathrm{m}^3$) & 100 & 1871--1970 & không \\
AirPassengers & \texttt{AirPassengers} [DL6] & hằng tháng & hành khách hàng không quốc tế (nghìn người) & 144 & 01/1949--12/1960 & $\log$; huấn luyện 120, kiểm định 24 \\
DAX (log giá) & \texttt{EuStockMarkets}, cột DAX [DL4] & ngày giao dịch & log giá đóng cửa & 1860 & $t=1,\dots,1860$ & $\log P_t$; huấn luyện 1500, kiểm định 360 \\
\bottomrule
\end{tabular}
\end{adjustbox}
\end{table}

\subsection{Xử lý dữ liệu}

\noidung{Tiền xử lý}
Dữ liệu thô được kiểm tra tự động trước khi biến đổi: đúng số quan sát như Bảng~\ref{tab:dulieu}, không thiếu giá trị, không trùng dòng, thời gian tăng dần, và mọi giá trị dương với ba bộ sẽ lấy logarit. Sau đó áp dụng các phép biến đổi ở cột cuối của bảng. Với Hồ Huron, biến xu thế $x=\text{năm}-1920$ đặt gốc ở giữa giai đoạn để hệ số chặn là mực nước ở khoảng giữa kỳ. Với Nile, ta thêm biến chỉ báo cho các năm từ 1899 trở đi để thảo luận về đổi mức. Hai bộ dùng để dự báo được chia theo thời gian (không chia ngẫu nhiên, mục 2.3.3): AirPassengers lấy 120 tháng đầu (01/1949--12/1958) huấn luyện, 24 tháng cuối (01/1959--12/1960) kiểm định; DAX (log giá) lấy 1500 ngày đầu huấn luyện, 360 ngày cuối kiểm định.

\noidung{Khám phá dữ liệu}
Ta nhìn nhanh từng chuỗi qua bảng thống kê mô tả (Bảng~\ref{tab:tab_eda_descriptive}) và đồ thị tổng quan (Hình~\ref{fig:fig_eda_HU_overview}--\ref{fig:fig_eda_DP_overview}). Các nhận xét sau chỉ nhằm định hướng cho các chương sau:
\begin{itemize}[leftmargin=*,itemsep=1pt]
\item \emph{Hồ Huron:} mực nước quanh $579$ ft (độ lệch chuẩn $1.32$ ft), có xu thế giảm; ACF giảm dần.
\item \emph{Linh miêu:} dao động theo chu kỳ rõ; ACF đổi dấu xen kẽ theo chu kỳ đó.
\item \emph{DAX (suất sinh lợi):} trung bình gần $0$ ($0.065\%$), độ lệch chuẩn $1.03\%$, đuôi nặng (kurtosis dư $6.28$), lệch âm ($-0.55$); biến động tụ thành cụm, ACF hầu như nằm trong dải $95\%$.
\item \emph{Nile:} lưu lượng trung bình giảm từ $1098$ (trước 1899) xuống $850$ (từ 1899), tức một lần đổi mức.
\item \emph{AirPassengers (log):} xu thế tăng và mùa vụ rõ; ACF giảm rất chậm.
\item \emph{DAX (log giá):} tăng dài hạn; ACF gần $1$ ở mọi độ trễ vẽ, PACF chỉ đáng kể ở độ trễ $1$.
\end{itemize}
Đây mới là quan sát bằng mắt; chuỗi nào dừng, chuỗi nào không, được kiểm chứng ở Chương 4 và 5.

\input{\tabdir/tab_eda_descriptive.tex}
\figx{fig_eda_HU_overview}{HU: mực nước hồ Huron với xu thế tuyến tính OLS, ACF và PACF mẫu}
\figx{fig_eda_LY_overview}{LY: $\log_{10}$ số linh miêu bị bẫy, ACF và PACF mẫu}
\figx{fig_eda_DR_overview}{DR: suất sinh lợi log ngày của DAX, ACF và PACF mẫu}
\figx{fig_eda_NI_overview}{NI: lưu lượng sông Nile (trung bình trước và sau 1899), ACF và PACF mẫu}
\figx{fig_eda_AP_overview}{AP: log số hành khách (huấn luyện \textbar{} kiểm định), ACF và PACF mẫu}
\figx{fig_eda_DP_overview}{DP: log giá DAX (huấn luyện \textbar{} kiểm định), ACF và PACF mẫu}

\subsection{Phương pháp kiểm chứng}

Mỗi nội dung ở Chương 2 được kiểm chứng bằng dữ liệu thực, bằng mô phỏng, hoặc cả hai (cột cuối Bảng~\ref{tab:danhmuc}). Bốn nguyên tắc sau dùng chung.

\noidung{Đối chiếu với phần mềm}
Các đại lượng có công thức rõ ràng (tự tương quan mẫu, PACF bằng đệ quy Durbin--Levinson, thống kê Ljung--Box, ước lượng Yule--Walker, trọng số $\psi_j$, ACF lý thuyết của ARMA) được nhóm tự cài đặt theo công thức ở Chương 2 rồi so với thư viện \texttt{statsmodels} [13]. Cách làm này giúp hiểu rõ công thức và kiểm tra cả phần mềm lẫn phần cài đặt; các phép so sánh được viết thành kiểm thử tự động với sai khác cho phép không quá $10^{-8}$. Cùng một đại lượng có thể cho số khác nhau tùy cách tính (chẳng hạn sai số chuẩn của hệ số ARIMA tính từ tích ngoài của gradient hay từ nghịch đảo ma trận thông tin quan sát), nên khi báo cáo nhóm luôn nêu cách đã dùng.

\noidung{Mô phỏng}
Có những tính chất không thấy được từ một chuỗi duy nhất, như phân phối lấy mẫu của ước lượng, độ phủ của khoảng tin cậy hay phân phối của thống kê kiểm định. Khi đó ta dùng \emph{mô phỏng Monte Carlo}: lấy một mô hình (đặt trước hoặc đã khớp trên dữ liệu thực) làm ``sự thật'', sinh $B$ chuỗi cùng độ dài, tính đại lượng quan tâm trên từng chuỗi rồi tổng hợp. Kết quả mô phỏng cũng có sai số: với một tần suất $\hat p$ ước lượng từ $B$ lần lặp độc lập, \emph{sai số Monte Carlo} (MCSE) là
\[
\mathrm{MCSE}(\hat p)\approx\sqrt{\frac{\hat p(1-\hat p)}{B}}\quad(\approx0.005\text{ khi }B=2000,\ \hat p\approx0.95),
\]
còn với một trung bình mô phỏng, MCSE là độ lệch chuẩn mẫu chia cho $\sqrt B$. Khi so giá trị mô phỏng với lý thuyết, ta tính chênh lệch theo đơn vị MCSE: lệch trong vài MCSE vẫn phù hợp. Số lần lặp thường là $B=2000$ (một vài mô phỏng nặng hơn dùng $B=1000$); mọi chuỗi ARMA mô phỏng bỏ $200$ giá trị đầu (đoạn khởi động) để gần phân phối dừng.

Các mô phỏng và mô hình ``sự thật'' tương ứng gồm: AR(1) với $\phi=0.6$, $n=100$ (4.1(c)); AR(1) với $(\hat\phi,\hat\sigma)$ của Hồ Huron, $n=98$ (4.2(a)); MA(2) với $\hat\theta$ của Hồ Huron, $n=98$ (4.2(d)); nhiễu iid $N(0,1)$, $n=200$ (4.3(a)); AR(2) với $\hat\phi$ của Linh miêu, $n=114$ (4.3(c)); bước ngẫu nhiên có trôi với $(\hat\delta,\hat\sigma)$ của DAX và các bước ngẫu nhiên độc lập (5.1(b)); nhiễu trắng và bước ngẫu nhiên cho sai phân (5.3(a)); bước ngẫu nhiên $n=200$ và AR(1) với $\phi=0.95,\,0.9,\,0.8$ cho kiểm định nghiệm đơn vị (5.3(b)); ARIMA$(0,1,1)$ với $\theta=-0.5$ và với tham số Nile cho khoảng dự báo (6.2(a)).

\noidung{Đánh giá dự báo}
Với hai bộ dự báo, mô hình chỉ được khớp trên phần huấn luyện rồi dự báo và so với phần kiểm định. Để kết quả không phụ thuộc vào đúng một điểm chia, nhóm còn đánh giá cuốn chiếu: khớp lại mô hình tại mỗi gốc, dự báo nhiều chân trời, tính sai số riêng cho từng chân trời (mục 2.3.3). Với Nile, gốc đầu tiên khớp trên 80 năm, chân trời tối đa 5 năm; với AirPassengers, trên 108 tháng, chân trời tối đa 12 tháng; dự báo một bước trên phần kiểm định (AirPassengers từ tháng 121, Nile từ năm 81) cũng khớp lại tại từng gốc. Các mốc chuẩn là trung bình, naive, naive mùa vụ và trôi; độ chính xác đo bằng MAE, RMSE, MAPE, MASE; khác biệt giữa hai phương pháp được kiểm tra bằng Diebold--Mariano.

\noidung{Tái lập}
Các mô phỏng dùng bộ sinh số giả ngẫu nhiên của \texttt{numpy} với hạt giống cố định \texttt{3179}; mỗi mô phỏng nhận một dòng số riêng sinh ra từ hạt giống đó, nên thêm hay bớt một mô phỏng không đổi kết quả của các mô phỏng khác. Phần mềm gồm Python 3.12.14 với \texttt{statsmodels} 0.15.0, \texttt{numpy} 2.5.5, \texttt{pandas} 3.1.1, \texttt{scipy} 1.18.1, \texttt{matplotlib} 3.11.2 trong một môi trường ảo riêng. Toàn bộ quy trình, từ tải dữ liệu đến sinh bảng và hình, chạy được bằng một lệnh (Phụ lục B).

\subsection{Tiêu chí đánh giá}

Mỗi nội dung kiểm chứng kết thúc bằng một nhận xét xếp loại theo ba mức, dùng thống nhất ở mọi chương và đặt ra trước khi xem kết quả (Bảng~\ref{tab:thangdo}).

\begin{table}[H]
\centering\footnotesize
\caption{Thang xếp loại mức độ phù hợp}\label{tab:thangdo}
\renewcommand{\arraystretch}{1.25}
\begin{tabularx}{\linewidth}{@{}lXX@{}}
\toprule
Mức & Đại lượng tính toán hoặc mô phỏng & Kết quả trên dữ liệu thực \\
\midrule
Phù hợp & Sai khác với công thức hoặc thư viện không quá $10^{-8}$; giá trị mô phỏng lệch giá trị lý thuyết không quá $3$ MCSE. & Kết luận khớp lý thuyết ở mức ý nghĩa 5\% (hoặc trong dải $\pm1.96/\sqrt n$). \\
Phù hợp một phần & Đúng về dạng hoặc chiều hướng nhưng sai khác lớn hơn dung sai trên. & Đúng với một số bộ dữ liệu, độ trễ hoặc chân trời nhưng không với tất cả. \\
Không phù hợp & Trái với dạng hoặc chiều hướng mà lý thuyết nêu. & Trái với kết luận của lý thuyết ở hầu hết các bộ áp dụng. \\
\bottomrule
\end{tabularx}
\end{table}

Sai khác nhỏ do cách tính hoặc sai số số học (hai thư viện làm tròn hay tối ưu hơi khác nhau) được nêu trong nhận xét và không bị tính là không phù hợp; với ước lượng bằng tối ưu hóa số (như cực đại hợp lý), ta so ở mức làm tròn của giá trị báo cáo thay vì $10^{-8}$. Khi kết quả phù hợp một phần hoặc không phù hợp, nhóm giải thích nguyên nhân dựa trên đặc điểm dữ liệu hoặc điều kiện áp dụng của lý thuyết.
